\documentclass[11pt]{article}
\usepackage[margin=0.9in]{geometry}
\usepackage{amsmath}
\usepackage{amssymb}
\usepackage{graphicx}
\usepackage{booktabs}
\usepackage{xcolor}
\usepackage{float}
\usepackage{verbatim}
\usepackage[font=small]{caption}
\graphicspath{%
  {../output/figures_sllg/track_width/hk36_box_comparison/}%
  {../output/figures_sllg/track_width/rkky_breaking/}%
  {../output/figures_sllg/track_width/hk36_D0p72_350x500/}%
  {../output/figures_sllg/track_width/hk36_D0p72_700x500/}%
  {../output/figures_sllg/track_width/hk36_D0p72_1400x500/}%
  {../output/figures_driving_T/}%
  {../output/figures_driving_T/inertia/}%
  {../output/figures_driving_T/campaign_dc/}%
  {../output/figures_driving_T/ranking/}%
  {../output/figures_driving_T/stability/}%
  {../output/figures_driving_T/width/}}

\newcommand{\Q}{\lvert Q\rvert}

\title{Finite-temperature skyrmion stability\\[2pt]
\large SAF racetrack, $H_{k,\mathrm{top}}=36$~mT, $D=0.72$~mJ/m$^2$,
$L_x\in\{700,1400,2800\}$~nm}
\author{}
\date{}

\begin{document}
\maketitle

\section{Setup}
Three racetracks differing only in length $L_x$ along the periodic
drive axis, at fixed width $L_y=1000$~nm (free $y$). Each $(T,J)$ cell
is an ensemble of $100$ independent thermal realizations. Every
realization is classified from its final field into compact skyrmion
(S), elongated (E), labyrinth (L) or annihilated (A);
$P_{\mathrm{surv}}=[n(\mathrm{S})+n(\mathrm{E})]/n_{\mathrm{ens}}$.

\begin{table}[H]
\centering
\caption{Material, geometry and grid.}
\begin{tabular}{lll}
\toprule
quantity & value & note\\
\midrule
$M_s$                 & $1.43\times10^{6}$~A/m & \\
$A_{\mathrm{ex}}$     & $16$~pJ/m & \\
$\alpha$              & $0.14$ & \\
$K_{\mathrm{top}}$    & $1.3106\times10^{6}$~J/m$^3$
                      & $H_{k,\mathrm{top}}=36$~mT\\
$K_{\mathrm{bot}}$    & $1.31\times10^{6}$~J/m$^3$ & \\
$D$                   & $0.72\times10^{-3}$~J/m$^2$ & \\
$t_{\mathrm{Co}}$     & $1.3$~nm & \\
$d_{\mathrm{Ru}}$     & $1.35$~nm & \\
$H_{\mathrm{RKKY}}$   & $0.205$~T & SAF interlayer coupling\\
cell $a$              & $2$~nm & \\
$\Delta t$            & $50$~fs & all three stages\\
$L_x$                 & $700 / 1400 / 2800$~nm
                      & $n_x = 350 / 700 / 1400$\\
$L_y$                 & $1000$~nm & $n_y = 500$, free $y$\\
$T_{\mathrm{sub}}$    & $10, 50, 100, 130, 160, 200$~K & 6 values\\
$J$                   & $0.5, 1, 2, 3, 4, 5\times10^{11}$~A/m$^2$
                      & 6 values\\
ensembles per cell    & $100$ & independent thermal seeds\\
trajectories per box  & $3600$
                      & $3594$ for $L_x=2800$~nm\\
demag / BC            & racetrack & periodic $x$, free $y$\\
\bottomrule
\end{tabular}
\end{table}

\begin{table}[H]
\centering
\caption{Three-stage protocol: durations, stopping rules and
thresholds.}
\begin{tabular}{lll}
\toprule
quantity & value & note\\
\midrule
\multicolumn{3}{l}{\emph{Stage 1 -- $T=0$ relaxation (shared seed
per box)}}\\
relaxation chunk      & $2000$ steps & $0.1$~ns per sample\\
number of chunks      & $125$ & no early stopping\\
total relaxation time & $250\,000$ steps $=12.5$~ns & fixed-step\\
convergence check     & size / energy plateau
                      & per-chunk, plus torque probe\\
\midrule
\multicolumn{3}{l}{\emph{Stage 2 -- thermal equilibration (per
$T_{\mathrm{sub}}$, per realization)}}\\
axis tolerance        & $0.02$ (relative) & on $D_1$, $D_2$\\
stable windows        & $3$ consecutive & stop condition\\
step cap              & $300\,000$ steps $=15$~ns
                      & reached only at high $T$\\
$\lvert m\rvert$ tolerance & $5\times10^{-3}$
                      & renormalization guard\\
demag convergence     & $0.02$ & image-sum tolerance\\
\midrule
\multicolumn{3}{l}{\emph{Stage 3 -- current drive}}\\
drive time            & $40\,000$ steps $=2$~ns
                      & no relaxation in this stage\\
sampling interval     & $200$ steps $=10$~ps & $200$ samples\\
field-dump interval   & $400$ steps $=20$~ps
                      & $103$ frames, ens0 only\\
survival threshold    & $\Q\ge0.5$ & topological charge\\
consecutive samples   & $10$ & sustained before a cell is called dead\\
\bottomrule
\end{tabular}
\end{table}

\noindent Stage 1 is run once per box and shared by all $600$
equilibration runs of that box; stage 2 caches one thermal state per
$(T_{\mathrm{sub}},\text{realization})$, which stage 3 then drives at
each $J$. The classifier is the one of the earlier report: topology
first ($\Q$), then domain count and solidity, then shape -- a
single-winding convex domain is a skyrmion, tagged \textbf{E} if it
percolates ($D_x/L_x\ge0.5$) or is elongated ($D_1/D_2\ge1.7$), else
\textbf{S}. The gate $D_x/L_x\ge0.5$ is the only ingredient that
refers to the box.

\section{Stability maps}

\begin{figure}[H]
\centering
\includegraphics[width=\linewidth]{stability_maps.png}
\caption{Ensemble stability regimes for the three track lengths.
Colour is the dominant class over $100$ realizations; the number in
each cell is $P_{\mathrm{surv}}$. The survival boundary
($T\simeq100$--$130$~K, and the high-$J$ edge) is at the same place in
all three boxes; what moves is the internal S/E boundary, which
retreats to higher $J$ and higher $T$ as the track is lengthened.}
\end{figure}

\begin{figure}[H]
\centering
\includegraphics[width=\linewidth]{composition_hk36_D0p72_350x500.png}
\caption{$L_x=700$~nm: per-class ensemble fractions. The annihilated
channel is not shown because it is identically zero here and in both
longer tracks: at this anisotropy the skyrmion is never lost by
collapse to the ferromagnet, only by stripe/labyrinth formation.}
\end{figure}

\begin{figure}[H]
\centering
\includegraphics[width=\linewidth]{composition_hk36_D0p72_700x500.png}
\caption{$L_x=1400$~nm: per-class ensemble fractions. Weight has moved
from E to S relative to the $700$~nm track, while the L map -- the
actual loss channel -- is unchanged.}
\end{figure}

\begin{figure}[H]
\centering
\includegraphics[width=\linewidth]{composition_hk36_D0p72_1400x500.png}
\caption{$L_x=2800$~nm: per-class ensemble fractions. Several cells are
mixed rather than pure -- $T=100$~K, $J=2\times10^{11}$ is
$0.88\,$S$/0.11\,$E, $J=3\times10^{11}$ is
$0.17\,$S$/0.76\,$E$/0.07\,$L, and $T=50$~K, $J=5\times10^{11}$ is
$0.65\,$E$/0.35\,$L. Comparing the three figures, the S and E maps
exchange weight as $L_x$ grows while the L map stays put.}
\end{figure}

\section{Final configurations}

\begin{figure}[H]
\centering
\includegraphics[width=\linewidth,height=0.86\textheight,%
  keepaspectratio]{config_table_hk36_D0p72_350x500.png}
\caption{$L_x=700$~nm: final driven top-layer $m_z$ of every $(T,J)$
cell (ens0 realization), drawn at physical scale on the stability
grid, labelled with the ens0 class; the suffix \textbf{p} marks cells
whose $x$-extent trips the periodic-image gate. Blue is the reversed
domain, red the background. Only the top layer is shown; the bottom
layer mirrors it under the SAF coupling. Almost the whole surviving
region carries a \textbf{p}: the domains fill the short track.}
\end{figure}

\begin{figure}[H]
\centering
\includegraphics[width=\linewidth,height=0.86\textheight,%
  keepaspectratio]{config_table_hk36_D0p72_700x500.png}
\caption{$L_x=1400$~nm, same construction. The low-$T$ cells now show
a compact blob with clear background on both sides, and the
\textbf{p} tags are gone from the $T\le100$~K block.}
\end{figure}

\begin{figure}[H]
\centering
\makebox[\linewidth][c]{%
  \includegraphics[width=1.09\linewidth,height=0.86\textheight,%
    keepaspectratio]{config_table_hk36_D0p72_1400x500.png}}
\caption{$L_x=2800$~nm, same construction. The surviving objects are
visibly identical in size to those of the $1400$~nm track while
occupying half the frame, which is the whole point of the comparison:
the skyrmion does not know how long the track is. The
$J\ge3\times10^{11}$ column shows the genuine current-driven
elongation, and the $T\ge130$~K rows the labyrinth state.}
\end{figure}

\section{Survival probability and Hall angle}

\begin{figure}[H]
\centering
\includegraphics[width=0.93\linewidth]{survival_hall_boxes.png}
\caption{Survival probability (top row), ensemble-mean Hall angle
$\bar\theta_H$ (middle) and its standard error
$\sigma_{\theta_H}/\sqrt{n}$ (bottom); columns are the three tracks.
Each row uses one colour scale across all boxes, so the maps may be
read against each other directly. All quantities are over the full
$100$-member ensemble. The three survival maps are the same map;
$\bar\theta_H$ grows with $J$ at low $T$, and its error bar is large
only in the near-dead $T=130$~K row, where few realizations survive to
define an angle.}
\end{figure}

\section{Velocity versus current}

\begin{figure}[H]
\centering
\includegraphics[width=\linewidth]{velocity_boxes.png}
\caption{Drift speed versus current density, ensemble mean $\pm$ SE,
for cells with $P_{\mathrm{surv}}\ge0.5$; common axis limits across
the three tracks. The curves and their mild sublinearity are the same
in all three boxes; only the extent of the surviving region differs
(the $T=100$~K branch reaches $J=4\times10^{11}$ in the shortest track
only).}
\end{figure}

\begin{figure}[H]
\centering
\includegraphics[width=\linewidth]{transport.png}
\caption{The same drift speed, and the Hall angle, with all three
boxes overlaid on one pair of axes (solid $700$, dashed $1400$, dotted
$2800$~nm). The curves superpose: transport is independent of track
length.}
\end{figure}

\section{Thermal equilibration}

\begin{figure}[H]
\centering
\includegraphics[width=\linewidth]{equil_D1D2_boxes.png}
\caption{Stage-2 thermal equilibration: ensemble-mean ellipse axes
$D_1(t)$ (top row) and $D_2(t)$ (bottom row) with $\pm1\sigma$ bands,
one colour per $T_{\mathrm{sub}}$; columns are the three tracks, all
on common axis limits. Each series is truncated to the window shared
by all $100$ members. At $T\le50$~K the axes stay flat at the $T=0$
value ($\approx190$~nm): the seed is already the thermal equilibrium.
From $T=100$~K upward the domain grows without arresting -- the
run-away to the labyrinth that the stability maps report as
\textbf{L}. The $700$~nm track is the exception that proves the rule:
its $T\ge160$~K curves flatten near $D_1\approx690$~nm because the
domain has run into the length of the box, while both longer tracks
continue to $\approx940$~nm.}
\end{figure}

\section{Box-size comparison}

\begin{figure}[H]
\centering
\includegraphics[width=\linewidth]{class_flips.png}
\caption{Top: cells whose dominant class changes between consecutive
track lengths ($6$ cells for $700\!\rightarrow\!1400$~nm, $1$ cell for
$1400\!\rightarrow\!2800$~nm; $7$ of $36$ cells change at all).
Bottom: the corresponding change in survival probability, which is
zero or within ensemble noise everywhere except the single cell
$T=100$~K, $J=4\times10^{11}$.}
\end{figure}

\begin{figure}[H]
\centering
\includegraphics[width=\linewidth]{gate_dx_lx.png}
\caption{The periodic-image gate metric $D_x/L_x$ (ens0); a star marks
the cells the gate sends to class \textbf{E}. In the low-$T$ block the
metric halves with every doubling of $L_x$ ($T=10$~K:
$0.25\rightarrow0.13\rightarrow0.06$), i.e.\ $D_x$ itself is fixed and
only its ratio to the box changes. Every S/E flip in the previous
figure is a cell whose star disappears here.}
\end{figure}

\begin{figure}[H]
\centering
\includegraphics[width=\linewidth]{box_invariance.png}
\caption{Left: $P_{\mathrm{surv}}(T)$ per current, all boxes overlaid.
Centre: aspect ratio $D_1/D_2$ versus $J$ per temperature, all boxes
overlaid. Right: absolute $x$-extent $D_x$ of the surviving cells,
with the $0.5L_x$ gate of each box drawn in grey. The first two
collapse onto single curves. The third shows the mechanism: the
$T=100$~K branch sits at $D_x\approx450$--$580$~nm, above the $350$~nm
gate of the shortest track but below the $700$ and $1400$~nm gates of
the longer ones -- one and the same object, classified differently.}
\end{figure}

\clearpage
\section{Interlayer coupling under drive}

\begin{figure}[H]
\centering
\includegraphics[width=\linewidth]{alignment_vs_T.png}
\caption{Interlayer alignment of the SAF pair at $j=0$, from the
stored two-layer state at the end of the thermal equilibration. Left:
$\langle m_{\rm top}\cdot m_{\rm bot}\rangle$ against the locked
($-1$) and unlocked ($0$) references. Right: the deviation
$1+\langle m_{\rm top}\cdot m_{\rm bot}\rangle$, which is linear in
$T$ at $3.3\times10^{-4}$/K with no upturn -- equipartition
transverse fluctuation, not a weakening of the coupling. The lock is
box-independent and never breaks in this range: at $200$~K it is still
good to $7\%$, well above the $\sim130$~K labyrinth transition, so the
lock outlives the skyrmion.}
\end{figure}

\begin{figure}[H]
\centering
\includegraphics[width=\linewidth]{alignment_vs_t_drive.png}
\caption{Alignment during the $2$~ns drive for the ens0 realization,
at the lowest and highest current of each temperature. Each trace
settles to a plateau set by $T$ and $J$; the current costs at most a
further $\sim5\%$ of alignment. The $T=100$~K, $J=5\times10^{11}$
trace is the one case that rings, a damped oscillation of period
$\approx0.28$~ns decaying over $\sim1$~ns.}
\end{figure}

\begin{figure}[H]
\centering
\includegraphics[width=\linewidth]{stretch_vs_current.png}
\caption{Current-driven interlayer stretch. Left: the steady offset
$d=|r_{\rm top}-r_{\rm bot}|$ against current, for the three track
lengths (line style) and three temperatures (colour); the curves
coincide to better than $2\%$, so the stretch is a local property of
the pair. Centre: the same offsets against the drift speed, with the
proportional-drag reference $d=v\tau$. Right: the implied lag
$\tau=d/v$.}
\end{figure}

\paragraph{Elastic, not decoupling.}
The offset reaches only $6.9$~nm at $10$~K and $8.4$~nm at $50$~K for
$J=5\times10^{11}$~A/m$^2$, against a skyrmion diameter of
$\approx190$~nm -- at most $4\%$ of the object. The two layers stay
locked and simply stretch against the RKKY coupling; they do not come
apart.

\paragraph{The linear-drag picture holds only at low drive.}
A rigid pair dragged through a linear spring would satisfy $d=v\tau$
with a single $\tau$. Instead $\tau$ rises from $\approx9$--$11$~ps at
$J\le2\times10^{11}$ to $19$--$21$~ps at $5\times10^{11}$, and the
$T=100$~K branch falls well below the $10$ and $50$~K ones, so the
points do not collapse onto one line. Past $J\approx3\times10^{11}$ the
skyrmion elongates ($D_1/D_2\rightarrow2.9$) and stops behaving as a
rigid disc, and the stretch then grows faster than the speed.

\begin{figure}[H]
\centering
\includegraphics[width=\linewidth]{separation_hk36_D0p72_350x500.png}\\[2pt]
\includegraphics[width=\linewidth]{separation_hk36_D0p72_700x500.png}
\caption{Separation dynamics for $L_x=700$~nm (top) and $1400$~nm
(bottom): the ensemble-mean $d(t)$ of every surviving cell (left) and
the steady offset against current (right). $d(t)$ rises and plateaus;
the detrended second half has an FFT peak-to-mean of only $2$--$4$
(a coherent oscillation over $100$ samples would give $\sim50$), so
apart from the single ringing cell above there is no sustained pair
oscillation.}
\end{figure}

\begin{figure}[H]
\centering
\includegraphics[width=\linewidth]{separation_hk36_D0p72_1400x500.png}
\caption{Separation dynamics for $L_x=2800$~nm, identical to the two
shorter tracks.}
\end{figure}

\begin{figure}[H]
\centering
\includegraphics[width=\linewidth]{pair_maps_hk36_D0p72_350x500.png}\\[2pt]
\includegraphics[width=\linewidth]{pair_maps_hk36_D0p72_700x500.png}\\[2pt]
\includegraphics[width=\linewidth]{pair_maps_hk36_D0p72_1400x500.png}
\caption{Pair coordinates at the drive end for $L_x=700$, $1400$ and
$2800$~nm (top to bottom): separation and charge imbalance
$|Q_{\rm top}+Q_{\rm bot}|$, both over all $100$ realizations and
masked to cells with $P_{\rm surv}\ge0.5$. Within the surviving region
the separation stays at the few-nm level and the imbalance below
$\sim0.01$. The large separations and imbalances of $0.15$--$0.22$ seen
outside this region belong to cells that are already labyrinths, where
a single pair coordinate has no meaning -- which is why they are
masked out.}
\end{figure}

\clearpage
\section{DC drive: velocity and deformation}
\label{sec:dcvdef}

Each curve
is an ensemble average over the $100$ realizations of a cell; the band
or bar is the standard error. Members are gated to single-winding,
\emph{$x$-localized} skyrmions: compact (S), elongated (E), and
spanning-flagged (P) members that are still localized along the
periodic axis ($D_x/L_x<0.5$). The P flag fires when the reversed
domain touches both $x$-edges; on a periodic track, that also catches a
skyrmion merely sitting on the seam. Only a domain that truly
spans the track ($D_x/L_x\ge0.5$), or a labyrinth (L), annihilated (A)
or reversed-background (R) member, is dropped. Each figure shows all
three tracks, top to bottom $L_x=700,1400,2800$~nm (panel titles label
them by cell count $n_x=350,700,1400$).

\begin{figure}[H]
\centering
\includegraphics[width=\linewidth]{S41_vt_hk36_D0p72_350x500.png}\\[2pt]
\includegraphics[width=\linewidth]{S41_vt_hk36_D0p72_700x500.png}\\[2pt]
\includegraphics[width=\linewidth]{S41_vt_hk36_D0p72_1400x500.png}
\caption{S41. Ensemble-mean drift speed $\langle|v|\rangle(t)$ under
the $2$~ns DC drive, one trace per peak current, at $T=10,50,100$~K,
for $L_x=700,1400,2800$~nm (top to bottom). A current is one colour
across panels; the $\pm$SE band is thinner than the line. The speed
rises to a current-set plateau within $\sim0.3$~ns and holds -- a
steady drift, not a transient. At $T=100$~K only the weakest currents
survive; the stronger drives break the skyrmion into a labyrinth and
are gated out. The $700$~nm track has no $T=100$~K survivors at all
(empty panel) -- its thermally swollen skyrmion already spans the
short track. The plateau speeds agree across the tracks that do
survive.}
\end{figure}

\begin{figure}[H]
\centering
\includegraphics[width=\linewidth]{S44_A_D1D2_t_hk36_D0p72_350x500.png}\\[2pt]
\includegraphics[width=\linewidth]{S44_A_D1D2_t_hk36_D0p72_700x500.png}\\[2pt]
\includegraphics[width=\linewidth]{S44_A_D1D2_t_hk36_D0p72_1400x500.png}
\caption{S44A. Ensemble-mean ellipse axes $D_1$ (major) and $D_2$
(minor) versus time, at the largest surviving current of each
temperature, for $L_x=700,1400,2800$~nm (top to bottom). The major
axis stretches monotonically while the minor axis stays nearly fixed:
the drive elongates the skyrmion, chiefly across the free-$y$ width,
and the stretch has not saturated by $2$~ns at the strongest drives.}
\end{figure}

\begin{figure}[H]
\centering
\includegraphics[width=\linewidth]{S44_B_vD_J_hk36_D0p72_350x500.png}\\[2pt]
\includegraphics[width=\linewidth]{S44_B_vD_J_hk36_D0p72_700x500.png}\\[2pt]
\includegraphics[width=\linewidth]{S44_B_vD_J_hk36_D0p72_1400x500.png}
\caption{S44B. Drive-average speed $v_{\mathrm{avg}}$ (left axis) and
the extremal axes $\max D_1$, $\min D_2$ (right axis) versus peak
current, per temperature, with SE bars, for $L_x=700,1400,2800$~nm
(top to bottom). $v_{\mathrm{avg}}$ rises sub-linearly with $J$; the
major axis grows with $J$ while the minor axis is nearly flat,
quantifying the drive-driven elongation.}
\end{figure}

\clearpage
\section{Pulse response -- setup}

The stability results above are for a DC drive: the current is switched
on and left on for 2~ns. This section sets up a separate study of the
\emph{pulsed} response, where the current has a finite duration and a
shape, and reports the usable-current window measured for each shape.

\paragraph{Geometry and material.}
Identical to the DC study: SAF racetrack, periodic $x$ and free $y$,
$H_{k,\mathrm{top}}=36$~mT (Set A), $D=0.72$~mJ/m$^2$, $a=2$~nm,
$\Delta t=50$~fs, on the same three tracks $L_x = 700$, $1400$ and
$2800$~nm at $L_y=1000$~nm. Each run is seeded from the DC campaign's
own cached states -- the $T=0$ relaxed configuration, or the
thermalised state of the matching $(T_{\mathrm{sub}},\text{realization})$
-- so the pulsed and DC results start from the same ensemble.

\paragraph{The six shapes.}
All are compared at matched peak amplitude $J_0$ and matched duration
$t_p=500$~ps, followed by a $500$~ps unforced interval so the skyrmion
is observed coming to rest. Delivered charge $\int J\,dt$ and action
$\int J^2\,dt$ then differ between shapes by construction, and are the
denominators of the efficiency measures.

\begin{table}[H]
\centering
\caption{Pulse shapes, with delivered charge and action at peak $J_0$
over duration $t_p$. The three triangles share both integrals
regardless of where the peak sits, so comparing them isolates asymmetry
alone.}
\begin{tabular}{llll}
\toprule
shape & profile & $\int J\,dt$ & $\int J^2\,dt$\\
\midrule
square & constant over $t_p$ & $J_0 t_p$ & $J_0^2 t_p$\\
half-sine & $\sin(\pi t/t_p)$ & $2J_0t_p/\pi$ & $J_0^2 t_p/2$\\
triangle, sharp rise & peak at $t=0$, linear fall
  & $J_0t_p/2$ & $J_0^2t_p/3$\\
triangle, sharp fall & linear rise, peak at $t_p$
  & $J_0t_p/2$ & $J_0^2t_p/3$\\
triangle, symmetric & peak at $t_p/2$
  & $J_0t_p/2$ & $J_0^2t_p/3$\\
Gaussian & centred, FWHM $250$~ps
  & $J_0\sigma\sqrt{2\pi}$ & $J_0^2\sigma\sqrt{\pi}$\\
\bottomrule
\end{tabular}
\end{table}

\begin{figure}[H]
\centering
\includegraphics[width=\linewidth,height=0.88\textheight,%
  keepaspectratio]{pulse_profiles.png}
\caption{The six drive profiles at matched peak amplitude and duration
(left), with their running delivered charge (centre) and running action
(right); the dotted line marks the end of the drive. Columns share
their vertical scale, so the shapes are directly comparable. The
endpoints are the values tabulated above: the square delivers twice the
charge of any triangle and three times the action, which is why it
destabilises the skyrmion at the lowest peak current. The three
triangles reach identical endpoints by different routes -- sharp-rise
front-loads and saturates by $\sim300$~ps, sharp-fall accumulates almost
nothing early and then climbs into the abrupt cut-off -- so comparing
them isolates the effect of rise/fall asymmetry at fixed charge and
fixed dissipation.}
\end{figure}

\paragraph{Usable-current window.}
A cap is the largest peak current at which \emph{every} realization
stayed a localized skyrmion, with every lower current also surviving
-- the contiguous window, not the maximum over surviving cells. Two
further classes are scored apart from the compact (S) and elongated (E)
skyrmions and count as failures: a spanning stripe (P) reaching its own
periodic image, and a reversed background (R) where the whole track has
switched polarity. Caps are quoted at finite temperature only, a $T=0$
point being a single deterministic realization.

\begin{table}[H]
\centering
\caption{Usable peak current ($10^{11}$~A/m$^2$) per pulse shape,
temperature and track length, from $2424$ pilot trajectories.
$^{*}$ marks a non-contiguous usable set. Every entry agrees across the
three lengths. The $700$~nm track has no $100$~K row: there the
thermally swollen skyrmion already spans $0.68$ of the track at the
lowest current tested, before the drive acts, so those cells measure
the box rather than the physics and are excluded ($1400$~nm sits at
$0.34$ under the same conditions).}
\label{tab:pulsecaps}
\begin{tabular}{lccccccccc}
\toprule
& \multicolumn{2}{c}{$L_x=700$~nm}
& \multicolumn{3}{c}{$L_x=1400$~nm}
& \multicolumn{3}{c}{$L_x=2800$~nm}\\
\cmidrule(lr){2-3}\cmidrule(lr){4-6}\cmidrule(lr){7-9}
shape & 10 & 50 & 10 & 50 & 100 & 10 & 50 & 100\\
\midrule
square              & 6 & 5 & 6 & 5 & 3 & 6 & 5 & 3\\
half-sine           & 8 & 6 & 8 & 6 & 4 & 8 & 6 & 4\\
triangle sharp rise & 9 & 6 & 9 & 6 & 5 & 9 & 6 & 5\\
\textbf{triangle sharp fall}
                    & \textbf{10} & \textbf{9} & \textbf{10}
                    & \textbf{9} & \textbf{6} & \textbf{10}
                    & \textbf{9} & \textbf{6}\\
triangle symmetric  & 9$^{*}$ & 6 & 9$^{*}$ & 6 & 5 & 9 & 6 & 5\\
Gaussian            & 9 & 6 & 9 & 6 & 5 & 9 & 6 & 5\\
\bottomrule
\end{tabular}
\end{table}

\section{Pulse response -- shape ranking}
\label{sec:shaperank}

With the usable window known, the question is which shape drives a
skyrmion best, according to three measures:
\begin{itemize}
\item displacement per delivered charge $\int J\,dt$; \\
\item displacement per delivered action $\int J^2\,dt$
(the Ohmic dissipation, i.e.\ energy); \\
\item and the largest displacement reached while the skyrmion stays compact.
\end{itemize}
The displacement is the net centroid travel, ensemble-averaged over the
surviving realizations of each cell; the charge and action are the
analytic integrals of the shape at its peak amplitude, so they are 
shape-specific. The three triangles share both integrals, so
any difference between them is pure rise/fall asymmetry. Members are
gated to $x$-localized skyrmions for the efficiency measures and to
compact (S) members for the reach measure. 
The short track carries no periodic wrap-around:
its largest displacement ($\sim\!218$~nm) is well below
$L_x=700$~nm for this case.

\subsection{Displacement}

\begin{figure}[H]
\centering
\includegraphics[width=\linewidth]{rank_d_per_charge_hk36_D0p72_350x500.png}\\[2pt]
\includegraphics[width=\linewidth]{rank_d_per_charge_hk36_D0p72_700x500.png}\\[2pt]
\includegraphics[width=\linewidth]{rank_d_per_charge_hk36_D0p72_1400x500.png}
\caption{Displacement per delivered charge, $d/\!\int J\,dt$, versus
peak current, one curve per shape, at each temperature, for
$L_x=700,1400,2800$~nm (top to bottom; all panels and boxes share the
$y$-range). This ratio corresponds to the skyrmion mobility: $v=\mu J$, so
$\mu = d/\!\int J\,dt$ (displacement per delivered charge, since
$d=\int v\,dt = \mu\!\int J\,dt$). At low drive all shapes collapse to
the linear-response mobility $\mu$ -- a property of the
skyrmion, not the pulse --, then separate as the square (blue) falls
fastest and the triangles stay highest. 
$\mu$ grows with temperature. 
The short track has no $T=100$~K survivors (empty panel).}
\end{figure}

\begin{figure}[H]
\centering
\includegraphics[width=\linewidth]{rank_d_per_action_hk36_D0p72_350x500.png}\\[2pt]
\includegraphics[width=\linewidth]{rank_d_per_action_hk36_D0p72_700x500.png}\\[2pt]
\includegraphics[width=\linewidth]{rank_d_per_action_hk36_D0p72_1400x500.png}
\caption{Displacement per delivered action (energy),
$d/\!\int J^2\,dt$, versus peak current, for $L_x=700,1400,2800$~nm
(top to bottom; shared $y$-range). The ordering is the same as for
charge but starker: the square is the least efficient at every current
and temperature, the triangles and Gaussian the most. The $1/J$
falloff is expected -- action grows as $J^2$ while
displacement grows as $J$. Triangles win because they produce
displacement while dissipating the least energy (the lowest
$\int J^2$): the square is the least energy-efficient shape and the
triangles the most.}
\end{figure}

\begin{figure}[H]
\centering
\includegraphics[width=0.8\linewidth]{rank_maxd_compact_hk36_D0p72_350x500.png}\\[2pt]
\includegraphics[width=0.8\linewidth]{rank_maxd_compact_hk36_D0p72_700x500.png}\\[2pt]
\includegraphics[width=0.8\linewidth]{rank_maxd_compact_hk36_D0p72_1400x500.png}
\caption{Largest displacement reached while the skyrmion stays compact
(S), per shape and temperature, for $L_x=700,1400,2800$~nm (top to
bottom; shared $y$-range). The ordering is the \emph{inverse} of the
efficiency one: the square moves a compact skyrmion the farthest
($\sim\!190$~nm) because it delivers the most charge, while the
efficient shapes trade reach for economy. The reach is consistent
across the three tracks.}
\end{figure}

\paragraph{Efficiency: triangular pulses preferred.}
At low drive every shape delivers the same displacement per charge,
the mobility $\mu$; at matched peak the square delivers twice a 
triangle's charge and three times its action, 
so it over-drives and deforms the skyrmion and its
efficiency falls fastest. Per unit dissipated energy the square is the
worst at every current. A triangle or a Gaussian is the choice to
move a skyrmion a given distance at least cost.

\paragraph{Reach: square pulses preferred.} 
The previous ordering inverts for absolute
reach: the square pushes a still-compact skyrmion the farthest,
since because it deposits the most charge. 

\paragraph{Asymmetry shows little effect.} The three triangles, matched in
charge and action, rank within a few percent on every metric.
Sharp-fall (linear rise into an abrupt cut-off) is marginally the best,
holding the largest compact reach at $T=100$ K, consistent with its
being the most stable shape in the cap table.

\paragraph{Displacement, side by side.} The trade-off is clearest in the
final configurations themselves. The next two figures place the worst
and a best shape, square (top) and the sharp-fall triangle (bottom), in
the same column, with a light-green arrow from the skyrmion's start to
its end (open circle at the end). At matched peak, the square moves the
skyrmion about twice as far (it delivers twice the charge) and deforms
it far more, elongating across the width; the triangle moves less but
stays compact. The two temperatures use different current ladders,
$J=2,4,6\times10^{11}$~A/m$^2$ at $T=10$~K and the lower $J=1,2,3$ at
$T=100$~K, where the stability window is narrower.

\begin{figure}[H]
\centering
\includegraphics[width=\linewidth]{dcompare_T010_hk36_D0p72_700x500.png}
\caption{Final top-layer $m_z$ (ens0) at $T=10$~K, square (top) vs
sharp-fall triangle (bottom), at $J=2,4,6$. The light-green arrow runs
from the start to the end (open circle at the end); the title gives the
net displacement. The square moves $\sim\!2\times$ farther and deforms
much more at every current.}
\end{figure}

\begin{figure}[H]
\centering
\includegraphics[width=\linewidth]{dcompare_T100_hk36_D0p72_700x500.png}
\caption{Same at $T=100$~K, at the lower $J=1,2,3$ where both shapes
stay within the narrower stability window. The skyrmion is thermally
swollen and the background noisier, but both shapes still displace
coherently; the square again moves farther per pulse.}
\end{figure}

\subsection{Speed}

\paragraph{Speed.} The same runs give the peak and average speed of the
skyrmion during the pulse. Both rise with peak current -- the peak
instantaneous speed almost linearly. Here the
square is the \emph{fastest}: sustaining $J_0$ for the whole pulse lets
the velocity reach its plateau, whereas a triangle holds $J_0$ only
briefly and the $\sim\!100$~ps inertial lag (next section) keeps the
skyrmion from ever reaching full speed. So the square wins on speed and
absolute reach, the triangles on efficiency.

\begin{figure}[H]
\centering
\includegraphics[width=\linewidth]{speed_max_hk36_D0p72_700x500.png}
\caption{Peak pulse speed versus peak current, one curve per shape, per
temperature, for the $L_x=1400$~nm track. The square (blue) reaches the
highest peak speed at every current; the triangles lag because their
current is at $J_0$ only briefly.}
\end{figure}

\begin{figure}[H]
\centering
\includegraphics[width=\linewidth]{speed_avg_hk36_D0p72_700x500.png}
\caption{Average pulse speed versus peak current, same layout, and on
the same $y$-scale as the peak speed above, so the two are directly
comparable. The square again leads, by a wider margin than for the
peak speed, because its current, and so its velocity, is sustained
rather than transient.}
\end{figure}

\subsection{Mobility}
The mobility $\mu = d/\!\int J\,dt$, where $v=\mu J$ in the linear-response 
regime,
is plotted against temperature below. It carries two clear trends. It
rises steeply with temperature, roughly $50\%$ from $T=0$ to $100$~K,
as thermal fluctuations help the skyrmion past the pinning and
deformation that otherwise limit its travel; and it falls with current:
at higher drive the skyrmion elongates and its displacement no
longer scales with the delivered charge. The spread over the six shapes
(the band width) is negligible at
$J=1\times10^{11}$~A/m$^2$, confirming that $\mu$ there does not depend 
on the pulse, and it widens only at
high current, where the square's mobility drops below the triangles' as
it over-drives.

\begin{figure}[H]
\centering
\includegraphics[width=0.72\linewidth]{mobility_range_hk36_D0p72_700x500.png}
\caption{Skyrmion mobility $\mu = d/\!\int J\,dt$ versus temperature on
the $L_x=1400$~nm track, one band per peak current. The band spans the
minimum to maximum over the six pulse shapes and the line is their
mean. $\mu$ rises with temperature (thermal assistance) and falls with
current (deformation). The band is narrow at low current, where $\mu$
is the shape-independent intrinsic mobility, and widens at high current
as the shapes diverge. $J=5\times10^{11}$ has no survivors above
$T=100$~K and $J=8\times10^{11}$ none above $T=50$~K.}
\end{figure}

\subsection{Operating points}
Combining efficiency and speed, each temperature has two extremal
operating points among the surviving cells: the most energy-efficient
(highest $d/\!\int J^2\,dt$) and the fastest (highest average pulse
speed). They sit at opposite corners of a Pareto front: the efficient
point is the weakest triangular pulse, it always corresponds to the
lower current density value, as less current is employed in deforming 
the skyrmion, and instead used to move it slowly;
whereas the fastest is a strong square pulse.

\begin{table}[H]
\centering
\caption{Energy-efficient and fastest operating points per temperature
for the $L_x=1400$~nm box, with the average pulse speed
$v_{\mathrm{avg}}$ (m/s) at each ($J$ in $10^{11}$~A/m$^2$). Efficiency
is maximized at the lowest current for a triangle; the fastest square
drops from $J=6$ to $J=3$ as the stability cap tightens with
temperature.}
\begin{tabular}{rllll}
\toprule
$T$ (K) & \multicolumn{2}{c}{energy-efficient}
& \multicolumn{2}{c}{fastest}\\
\cmidrule(lr){2-3}\cmidrule(lr){4-5}
 & shape ($J$) & $v_{\mathrm{avg}}$ & shape ($J$) & $v_{\mathrm{avg}}$\\
\midrule
0 & tri\_symmetric ($0.5$) & $23$ & square ($6.0$) & $361$\\
10 & tri\_sharpfall ($0.5$) & $28$ & square ($6.0$) & $366$\\
50 & tri\_sharprise ($0.5$) & $48$ & square ($5.0$) & $364$\\
100 & tri\_symmetric ($0.5$) & $58$ & square ($3.0$) & $315$\\
\bottomrule
\end{tabular}
\end{table}

\begin{figure}[H]
\centering
\includegraphics[width=0.72\linewidth]{operating_points_hk36_D0p72_700x500.png}
\caption{Energy efficiency $d/\!\int J^2\,dt$ versus average pulse
speed for every surviving cell of the $L_x=1400$~nm box, one colour per
temperature. Stars mark the most efficient operating point of each
temperature, diamonds the fastest. The two objectives trade off along
a Pareto front that temperature lifts upward, as the mobility grows.}
\end{figure}

\section{Pulse response -- duration}
\label{sec:duration}

The width sweep holds the peak current at each shape's cap and varies
the pulse duration $t_{\mathrm{pulse}}$ over $\{100,200,300,400,600,
700\}$~ps ($500$~ps is the production pulse), with $25$ realizations
per finite-$T$ cell. Since a longer pulse at fixed peak delivers
proportionally more charge, duration is a second drive knob. All
quantities are ensemble means over the $x$-localized survivors, on the
$L_x=1400$~nm box.

\begin{figure}[H]
\centering
\includegraphics[width=\linewidth]{width_displacement_hk36_D0p72_700x500.png}
\caption{Net displacement versus pulse duration, one curve per shape,
per temperature. Displacement grows almost linearly
with duration -- longer pulse, more delivered charge.
Curves stop where the cell no longer survives.}
\end{figure}

\begin{figure}[H]
\centering
\includegraphics[width=\linewidth]{width_speed_hk36_D0p72_700x500.png}
\caption{Average pulse speed versus duration. It rises and saturates:
a short pulse is dominated by the $\sim\!100$~ps inertial rise (next
section), which is a shrinking fraction of a longer pulse, so the
average approaches the plateau speed by $\sim\!400$--$600$~ps.}
\end{figure}

\begin{figure}[H]
\centering
\includegraphics[width=\linewidth]{width_deformation_hk36_D0p72_700x500.png}
\caption{Maximum ellipse major axis $D_1$ versus duration. The
skyrmion elongates roughly linearly with duration; at $T=100$~K it
starts thermally swollen and approaches the $\sim\!700$~nm track-half
(spanning) at the longest pulses, which is where survival falls.}
\end{figure}

\begin{figure}[H]
\centering
\includegraphics[width=\linewidth]{width_survival_hk36_D0p72_700x500.png}
\caption{Surviving-skyrmion fraction versus duration. Curves not shown
sit at $P_{\mathrm{surv}}=1$. At the cap current, extending the pulse
past $\sim\!400$--$600$~ps destabilizes the Gaussian, sharp-rise and
sharp-fall shapes, while the square, half-sine and symmetric triangle
survive every duration.}
\end{figure}

\paragraph{Duration is a stability axis.} The cap current was
calibrated at $t_{\mathrm{pulse}}=500$~ps, so a longer pulse at the
same peak over-drives the skyrmion. 
Notably the sharp-fall triangle, the most stable shape
at $500$~ps, is among the first to fail at $600$--$700$~ps.
Duration, peak current and shape are not
independent.

\begin{comment}
\paragraph{Energy partition (deferred).} The width runs also carry an
online translation-versus-deformation dissipation observable. It is
reliable only at $T=0$, where the fitted collective velocity matches
the tracked drift (there the square spends $\sim\!15$--$19\%$ of its
dissipation on translation and the rest on deformation); at finite
temperature the observable evaluates the deterministic $\mathrm{d}
\mathbf{m}/\mathrm{d}t$ on the thermally-roughened state, so the
incoherent roughness swamps the coherent drift and the partition
collapses. Its temperature dependence therefore awaits a corrected
re-run and is not reported here.
\end{comment}

\section{Pulse response -- inertia}
\label{sec:inertia}

The finite duration of a pulse exposes a transient the DC drive cannot:
the skyrmion accelerating when the current switches on and decelerating
when it switches off. We characterise it on the square pulse -- the
one shape with sharp, well-defined edges -- on the $L_x=1400$~nm track
($700\times500$ cells). The velocity is the top-layer drift speed
$|v|(t)=|\mathrm{d}\mathbf{r}/\mathrm{d}t|$ from the unwrapped centroid,
sampled every $10$~ps; the appendix gives the extraction of the
relaxation time in full.

\begin{figure}[H]
\centering
\includegraphics[width=0.82\linewidth]{S48_A_vt_full.png}
\caption{Complete $|v|(t)$ at $T=0$ for the square pulse, one trace per
peak current; the $500$~ps drive window is shaded. The velocity rises
to a plateau proportional to $J$, then decays back to a small residual
drift once the current is off. The response is \emph{overdamped}: no
overshoot at the plateau and no ringing after switch-off.}
\end{figure}

\begin{figure}[H]
\centering
\includegraphics[width=0.49\linewidth]{S48_B_vt_rise.png}
\hfill
\includegraphics[width=0.49\linewidth]{S48_C_vt_fall.png}
\caption{Rising (left) and falling (right) edges of $|v|(t)$, zoomed,
one trace per peak current at $T=0$. Both edges are single
exponentials to good approximation; the rise develops a mild overshoot
above the plateau only at the two highest currents.}
\end{figure}

\begin{figure}[H]
\centering
\includegraphics[width=0.49\linewidth]{S48_D_invtau.png}
\hfill
\includegraphics[width=0.49\linewidth]{S48_E_invtau_T.png}
\caption{Inverse relaxation time. Left: $1/\tau$ of the rise and fall
fits versus peak current at $T=0$; both sit at $\approx10$/ns and vary
by under $3\%$ across the whole current range, so $\tau\approx100$ ps is
essentially drive-independent, with $\tau_{\mathrm{rise}}\approx
\tau_{\mathrm{fall}}$ (symmetric turn-on/turn-off). Right: $1/\tau$
pooled over the rise and fall edges of every surviving ensemble member
versus temperature, at the common current $J=3\times10^{11}$ A/m$^2$;
the error bar is the ensemble standard deviation. It is flat within
scatter ($10.07,10.03,9.91,9.95$/ns at $T=0,10,50,100$ K), i.e., the
effective damping is temperature-independent.}
\end{figure}

The picture is a single overdamped relaxation time $\tau\approx100$~ps
($1/\tau\approx10$/ns), independent of drive current and of
temperature, and equal on turn-on and turn-off. There is no inertial
oscillation. Here the coupling is fixed and the response is overdamped. 

\begin{comment}
\clearpage
\section{Discussion}

\paragraph{What does not depend on the track length.}
The survival map, the morphology and the transport are box-invariant.
$P_{\mathrm{surv}}$ agrees cell-by-cell across the three tracks with a
mean absolute difference of $0.022$ ($700\!\rightarrow\!1400$~nm) and
$0.005$ ($1400\!\rightarrow\!2800$~nm), i.e.\ at the level of the
$1/\sqrt{100}=0.1$ ensemble noise of a single cell. The aspect ratio
$D_1/D_2$ reproduces to $\pm0.03$, the drift speed and the Hall angle
superpose within their standard errors, and the equilibration
transients at each $T_{\mathrm{sub}}$ are the same curves. The
physical boundaries of the phase diagram -- the thermal ceiling at
$T\simeq100$--$130$~K and the high-current edge -- are therefore
intrinsic to the material and the drive, not to the simulation cell.

\paragraph{What does depend on it.}
Only the S/E label, and only through the periodic-image gate. Seven of
$36$ cells change dominant class across the three boxes, all of them
in the surviving region, and all in the same direction: E at short
$L_x$, S at long $L_x$. The gate figure identifies the cause without
ambiguity -- $D_x$ is a fixed physical length (about $175$~nm at
$T=10$~K, $280$~nm at $T=50$~K, $450$--$580$~nm at $T=100$~K), so
$D_x/L_x$ halves each time the track is doubled and eventually drops
below the $0.5$ threshold. A skyrmion occupying $71\%$ of a $700$~nm
track and $21\%$ of a $2800$~nm one is the same skyrmion; calling the
first "elongated/spanning" is a statement about the box. The
$D_1/D_2\ge1.7$ branch of the classifier, which is box-free, flags the
genuinely stretched states ($J\ge3\times10^{11}$) identically in all
three boxes.

\paragraph{A second, milder finite-size effect.}
Above the stability ceiling the growing labyrinth in the $700$~nm
track runs into the box: its $D_1$ flattens at $\approx690$~nm, one
track length, while the $1400$ and $2800$~nm tracks agree with each
other at $\approx940$~nm. This does not affect any reported class --
those cells are \textbf{L} in every box -- but it means the short
track cannot be used to measure the size of the run-away state, only
to detect that it happened.

\paragraph{Convergence.}
The correction is not linear in $L_x$: six cells change between $700$
and $1400$~nm but only one between $1400$ and $2800$~nm. The shortest
track is qualitatively misleading -- it reports the entire $T=100$~K
row and half the $T=50$~K row as elongated -- whereas $L_x=1400$~nm
is already within one cell of the converged answer. For this material
$L_x\gtrsim1400$~nm ($\gtrsim5\times$ the largest $D_x$ of a surviving
skyrmion) is sufficient; $2800$~nm is converged.

\paragraph{The one physical box effect.}
The single cell where survival itself moves is $T=100$~K,
$J=4\times10^{11}$: $P_{\mathrm{surv}}=0.50$ at $L_x=700$~nm and
$0.00$ at $1400$ and $2800$~nm. This cell sits exactly on the burst
boundary, and the short track stabilizes it: a domain that has already
grown to $D_x\simeq0.72L_x$ closes on its own periodic image and is
held as a wrapped stripe, whereas in a long track the same domain has
room to keep elongating and breaks up into a labyrinth. The artifact
here is not a mislabelling but a spurious stabilization, and it acts
in the non-conservative direction -- the short box over-reports
survival. It is confined to this one boundary cell; the two longer
tracks agree.

\paragraph{Consequence.}
Quoting a stability map from a $700$~nm track would misreport the
skyrmion/elongated boundary over a third of the surviving region and
over-report survival at the burst edge. The $2800$~nm results are the
ones to quote; the $1400$~nm results are usable. Since $P_{\mathrm{surv}}$,
$v(J)$ and $\bar\theta_H$ are box-invariant, any earlier conclusion
resting on those quantities alone is unaffected by track length.
\end{comment}

\clearpage
\appendix
\section{Extraction of the velocity relaxation time}
\label{app:invtau}

The velocity is the magnitude of the time derivative of the unwrapped
top-layer centroid $\mathbf{r}(t)=(c_x,c_y)$,
\begin{equation}
|v|(t) = \left|\frac{\mathrm{d}\mathbf{r}}{\mathrm{d}t}\right|,
\end{equation}
evaluated by central differences on the $10$~ps sample grid. Writing
$t_p$ for the pulse duration ($500$~ps, read from each trajectory's
own configuration record), the two edges are fit to single
exponentials:
\begin{align}
\text{rise } (0\le t\le t_p):\quad
  &|v|(t) = A\left(1 - e^{-t/\tau_{\mathrm{rise}}}\right),
  \label{eq:rise}\\[2pt]
\text{fall } (t\ge t_p):\quad
  &|v|(t) = v_{\mathrm{res}}
    + (v_p - v_{\mathrm{res}})\,e^{-(t-t_p)/\tau_{\mathrm{fall}}},
  \label{eq:fall}
\end{align}
where $A$ is the plateau speed, $v_p$ the speed at switch-off and
$v_{\mathrm{res}}$ the residual drift. Each fit is a nonlinear
least-squares minimisation of $\sum_i\big(|v|_i-\text{model}\big)^2$
(Levenberg--Marquardt), and the reported quantity is the inverse time
constant $1/\tau_{\mathrm{rise}}$ or $1/\tau_{\mathrm{fall}}$.

The fall is fit over the full $[\,t_p,\,t_p+250\,\text{ps}\,]$ window.
The rise is fit over $[\,0,\,90\,\text{ps}\,]$, ending before the
mid overshoot; otherwise, a wider window would drag $A$ and
$\tau_{\mathrm{rise}}$ off their true values and inject a spurious
current dependence into $1/\tau_{\mathrm{rise}}$. The two windows and
the resulting fits are shown in Fig.~\ref{fig:invtaufit}: inside the
rise window the fit tracks the data, and its dashed extrapolation
climbs away from the plateau.

\begin{figure}[H]
\centering
\includegraphics[width=0.9\linewidth]{invtau_fit_windows.png}
\caption{The $1/\tau$ extraction on one representative cell (square
pulse, $J=5\times10^{11}$~A/m$^2$, $T=0$). Grey points are the measured
$|v|(t)$ (ground truth). Shaded bands are the rise (blue) and fall
(red) fit windows; solid curves are the exponential fits of
Eqs.~\eqref{eq:rise}--\eqref{eq:fall} inside their windows, dashed
where extrapolated. The dashed blue curve stops at $90$~ps to avoid 
the rising observable past $\sim\!110$~ps.}
\label{fig:invtaufit}
\end{figure}

For the temperature dependence (the $1/\tau$-versus-$T$ panel of
Sec.~\ref{sec:inertia}) the two values
$\{1/\tau_{\mathrm{rise}},1/\tau_{\mathrm{fall}}\}$ are collected from
every ensemble member that ends as a compact or elongated skyrmion --
the field classifier excludes labyrinth and reversed-background
members -- pooled at the common current $J=3\times10^{11}$~A/m$^2$; the
plotted point is their mean and the error bar their standard deviation.

\clearpage
\section{Per-shape stability of the pulse study}
\label{app:pulsestability}

For completeness, the two figures below classify every production
realization of the pulse-shape study over the $(T,J)$ grid, per shape,
for the $L_x=1400$~nm ($700\times500$) box. Since production is
capped at each shape's stability window and its current ladder stops
at $8\times10^{11}$~A/m$^2$, every cell in the window is a
compact (S) or, at the highest currents, an elongated (E) skyrmion at
$P_{\mathrm{surv}}=1.00$; blank cells lie beyond that shape's cap. The
other two tracks ($L_x=700$ and $2800$~nm) give the same picture and
are not reproduced here.

\begin{figure}[H]
\centering
\includegraphics[width=\linewidth]{pulse_stability_maps_hk36_D0p72_700x500.png}
\caption{Dominant final class in each $(T,J)$ cell, one panel per
shape, annotated with the surviving-skyrmion fraction
$P_{\mathrm{surv}}=(S+E+P)/n$. The legend lists all six classes
(compact S, elongated E, spanning P, labyrinth L, annihilated A,
reversed R); only S and E occur within the capped production window.
Blank cells are beyond the shape's cap.}
\end{figure}

\begin{figure}[H]
\centering
\includegraphics[width=\linewidth]{pulse_composition_hk36_D0p72_700x500.png}
\caption{Per-class ensemble fractions, rows the shapes and columns the
classes S, E, P and L. The weight sits almost entirely in S, shifting
to E only at the highest surviving currents; the P and L columns are
empty because production never enters those regimes.}
\end{figure}

\clearpage
\section{Open questions and operating guidance}
\label{app:questions}

Three practical questions about running a skyrmion racetrack at finite
temperature motivate the study; the results above bear on each.

\paragraph{Can the skyrmion keep the same speed at room temperature, or
must it be driven more slowly?} Within the surviving window the
mobility $\mu = d/\!\int J\,dt$ actually \emph{rises} with temperature
(Sec.~\ref{sec:shaperank}): from $\approx100$ to $\approx148$~nm per
$10^{11}$~A/m$^2$/ns between $T=0$ and $100$~K, so at fixed delivered
charge the skyrmion moves \emph{faster}, not slower, as it warms, since
thermal fluctuations help it past the pinning and deformation that
limit its travel. The caveat is survival: the compact skyrmion exists
only up to $T\simeq100$--$130$~K, above which it breaks into a
labyrinth, so true room temperature ($300$~K) is out of reach at this
anisotropy. Up to the stability ceiling, then, no slowdown is needed.

\paragraph{Is directional control harder with temperature?} On average
the drift stays along the current axis, but thermal fluctuations add a
transverse jitter: the ensemble spread in the skyrmion Hall angle
(Sec.~4, the survival/Hall maps) widens with temperature, so a single
realization deviates more from the mean deflection as it warms.
Directional control is therefore progressively harder toward the
labyrinth ceiling, even where the mean direction is unchanged.

\paragraph{What is the optimal drive temperature without annihilation?}
Because the mobility increases with temperature while the skyrmion
survives only up to $T\simeq100$--$130$~K, the best operating point is
just below that labyrinth transition, near $T\simeq100$~K for this
material: there the skyrmion is still compact
($P_{\mathrm{surv}}\approx1$ at low-to-moderate current) yet moves with
the highest mobility. Below it the skyrmion is stable but slower per
unit charge; above it, it is lost.

\end{document}